\documentclass[10pt,a4paper]{article}

\usepackage[margin=1.05in]{geometry}
\usepackage{setspace}
\usepackage{amsmath,amssymb,amsthm}
\usepackage{graphicx}
\usepackage{array}
\usepackage{caption}
\usepackage{enumitem}
\usepackage{titlesec}
\usepackage{microtype}
\usepackage{hyperref}
\usepackage{fancyhdr}

\hypersetup{colorlinks=true,linkcolor=black,urlcolor=black,citecolor=black}
\captionsetup{font=small,labelfont=bf,skip=4pt,justification=raggedright}

\titleformat{\section}{\normalsize\bfseries}{}{0em}{}
\titleformat{\subsection}{\small\bfseries}{}{0em}{}
\titleformat{\subsubsection}{\small\itshape}{}{0em}{}
\titlespacing*{\section}{0pt}{1.35em}{0.5em}
\titlespacing*{\subsection}{0pt}{0.95em}{0.35em}

\theoremstyle{plain}
\newtheorem{lemma}{Lemma}
\newtheorem{theorem}{Theorem}
\theoremstyle{definition}
\newtheorem{definition}{Definition}

\setstretch{1.14}
\setlist[itemize]{leftmargin=1.15em,itemsep=0.1em,topsep=0.2em}

\graphicspath{{drawings/}}

\pagestyle{fancy}
\fancyhf{}
\fancyhead[L]{\small\itshape Elements of Position}
\fancyhead[R]{\small\itshape VI.\ The Elevation}
\fancyfoot[C]{\small\thepage}
\renewcommand{\headrulewidth}{0pt}

\begin{document}

\begin{center}
\vspace*{2.2em}
{\large Elements of Position}\\[1.6em]
{\LARGE\bfseries VI.\ The Elevation}\\[2.2em]
{\normalsize Justin Erholtz}\\[0.35em]
{\small 15 September 2026}
\end{center}

\vspace{2.4em}

\begin{center}
{\large From outside the closed object, looking at it.}\\[0.55em]
{\large A meeting is a derived unit from one HERE.}
\end{center}

\vspace{1.8em}

\noindent
This part does not replace I--V.
It writes the third drawing.

Plan was the count, the pair $\{n,1/n\}$, the bags of first-writings.
Section was the local circle of radius $n/2$ and the half-walk $n\pi/2$.
Elevation is look-back from a pause: the weight of a finished tick seen from outside the unit it already closed.

IV asked what line a seated meeting may occupy.
The answer was the cut.
This part asks what formula, from one HERE, is allowed to vanish on that line.
The answer is not a new cut.
It is the dual sum of the complementary pair, wrapped by the local circle so the pause keeps one centre.

\vspace{1.4em}

\noindent
\textit{A note.}
Claims already proved in I--V are used, not reopened.
The drawings are layers on the same architectural set, built by \texttt{Elements of Position - Drawings.py}.

\vspace{0.9em}
\noindent
\textit{A note on names used here.}
Names from I--V stand.
Elevation --- look-back from a pause against what has already arrived.
Dated --- owned at a finished first-writing station, nothing later present.
Dual sum --- the two complementary interiors added.
Derived unit --- the reading of plan, section, and elevation together in the limit of finished stations.
Meeting of the derived unit --- a vanishing of a formula the stock can set to zero, from one HERE.

\newpage
\tableofcontents
\newpage

% =====================================================================
\section{What is already on the page}
% =====================================================================

I closed the unit and inverted the pair.
The lock is $1$.
The only primary self-ratio of a finished segment is $1/2$.
Rate is a title block.
Infinity is not a tick of the $n$-count.

II named $2$ from that $1$, kept the rate at $1$, and laid the other walk on the same midline.

III sat the meetings of those walks on that line.
The first exact marriage is at five.
Approximate dockings approach it.
The line does not move.

IV pulled the primary ratio back onto the index.
Let $\sigma=\tfrac12+\alpha$.
The complementary interiors at a finished arrival $n>1$ are
\[
\{n^{-\sigma},\,n^{-(1-\sigma)}\}.
\]
Their geometric mean is $n^{-1/2}$ for every real $\sigma$.
The ends are equal if and only if $\sigma=\tfrac12$.
Under rate $1$ the only permanent analytic cut of the derived unit is $\operatorname{Re}s=\tfrac12$.

A pause is one own-origin.
The ratio $\chi$ of local-circle samples at $1-s$ and at $s$ satisfies
\[
\bigl|\chi(\sigma+it)\bigr|=1
\quad\text{if and only if}\quad
\sigma=\tfrac12.
\]
Off the lock, $|\chi|$ is two radii.
The $n$-count does not seat two centres at once.

V pointed at $1$.
$1$ is THERE.
We are HERE.
Arrive at THERE and you are HERE.
The endless is already the interior of the unit.
Infinity is the other pole of that unit, not a later integer.

The missing drawing was the view from outside.

% =====================================================================
\section{The three drawings of one object}
% =====================================================================

The closed unit already owns three drawings.

\begin{itemize}
\item Plan: the axis, the $n$-count, the pair $\{n,1/n\}$, the nested packages, the first-writings and their bags.
\item Section: the local circle of radius $n/2$, the half-walk $n\pi/2$, the first lock $\pi$ at $n=2$.
\item Elevation: inversion through $1$, look-back against what has already been built.
\end{itemize}

None of these is the object.
The object is what all three are cutting.
A derived unit is the reading of the three together in the limit of finished stations.
A meeting of that unit is a vanishing from one HERE.

\begin{figure}[ht]
\centering
\includegraphics[width=0.95\textwidth]{layer_e01_three_drawings.png}
\caption{Three drawings of one object.
Plan keeps the face and the bags.
Section keeps the circle and the half-walk.
Elevation keeps the look-back.
The cut is common.}
\end{figure}

A frame frozen and flattened is a picture with no universal scale except its relation of parts.
Zoom is crop.
If there is a frame, there is a half-line.
Those three sentences were already forced.
Elevation is what you do with them when you stand outside the closed unit and ask the count to report.

% =====================================================================
\section{Look-back is logarithm}
% =====================================================================

IV already said it.
Look-back from a pause is logarithm.
Live ticks are additive.
The invariant of inversion is $dx/x$; its integral along a finished arrival $n>1$ is $\log n$.

That is the elevation of one tick.
The weight the tick contributes, seen from outside, is
\begin{equation}
\label{eq:weight}
n^{-s}=e^{-s\log n}.
\end{equation}
The two-part index $s=\sigma+it$ is the knife from IV: seen along the cut, unseen off it.
The unseen part turns the look-back.
It does not mint a second cut.

\begin{figure}[ht]
\centering
\includegraphics[width=0.92\textwidth]{layer_e02_lookback.png}
\caption{Look-back from a pause.
Live ticks sit on the axis.
The height of each look-back is $\log n$.
Tempo is not in the drawing.}
\end{figure}

Sum the weights.
That is the Dirichlet series of the count:
\begin{equation}
\label{eq:dirichlet}
Z(s)=\sum_{n=1}^{\infty}n^{-s}.
\end{equation}
The sum arrives at a finite mark only while $\sigma>1$.
At $\sigma=1$, $t=0$, the walk is the harmonic.
It does not close.
That is not a dead object.
It is a trace that will not arrive.

The same walk, regrouped by unique bags of first-writings, is the product
\begin{equation}
\label{eq:euler}
Z(s)=\prod_{p}\bigl(1-p^{-s}\bigr)^{-1},
\end{equation}
wherever the sum settles.
That sentence was already in the room: a sum over every frame-index is a product over the floors that will not reduce.
Plan wrote the bags.
Elevation writes the weights.
They are two drawings of one walk.

% =====================================================================
\section{Complementary interiors, written as a series}
% =====================================================================

Lemma~1 of IV, with the unseen part restored, is the pair
\[
n^{-s}+n^{-(1-s)}.
\]
On the cut $\sigma=\tfrac12$ the pair is real:
\begin{equation}
\label{eq:dual}
n^{-1/2-it}+n^{-1/2+it}
=2\,n^{-1/2}\cos(t\log n).
\end{equation}
The dual sum of the street is the elevation zeroed on the half,
\begin{equation}
\label{eq:dualsum}
D(t)=\sum_{n=1}^{\infty}2\,n^{-1/2}\cos(t\log n).
\end{equation}
Off the cut the two ends are unequal.
Their geometric mean is still $n^{-1/2}$.
Equality of ends is still only $\sigma=\tfrac12$.
Two leftovers are not two locks.

\begin{lemma}[Weights, with unseen]
At a finished arrival $n>1$ the pair $\{n^{-s},n^{-(1-s)}\}$ has geometric mean $n^{-1/2}$.
The two ends are equal if and only if $\sigma=\tfrac12$.
\end{lemma}

This is the first page of I, applied to the elevation.

% =====================================================================
\section{What a station owns}
% =====================================================================

A finished first-writing $p$ is a station.
The data permitted there are the arrivals $q\le p$, the local circle of radius $p/2$, and look-back against what has already arrived.
A station is not a meeting.
It is a place at which one may pause.

On that station the dated product
\begin{equation}
\label{eq:dated}
Z_p(s)=\prod_{q\le p}\bigl(1-q^{-s}\bigr)^{-1}
\end{equation}
is a finite object.
No arrival beyond $p$ is present.
Each factor $(1-q^{-s})^{-1}$ is never zero wherever it is defined.
So the dated product never vanishes.

\begin{lemma}[Dated product]
At every finished first-writing station $p$, $Z_p(s)$ has no zero.
\end{lemma}

The local-circle sample has no zeros in $0<\operatorname{Re}s<1$.
A midline polynomial $\tfrac12 s(s-1)$ accounts for the endpoints $0$ and $1$, and stops there.

Therefore a non-trivial vanishing is not a finite-depth fact.
It appears only when stations deepen, and only in a formula that is allowed to be zero.

The dated dual
\[
D_N(t)=2\sum_{n=1}^{N}n^{-1/2}\cos(t\log n)
\]
is the elevation at finite depth $N$.
Its crossings in $t$ migrate as $N$ grows.
They are crop.
A seated meeting would not pack up when the next first-writing arrives.

\begin{figure}[ht]
\centering
\includegraphics[width=0.92\textwidth]{layer_e04_dated_crop.png}
\caption{Dated dual on the cut.
Station $11$ and station $97$.
The crossings move.
Crop is not a seat.}
\end{figure}

On the side where the walk may arrive, crop dies.
$Z_p(2)$ approaches $\pi^2/6$.
That is the half-walk ratio squared over freeze $6$.
The $96$-gon of the elevation.
Genetic.
No finished inventory.

% =====================================================================
\section{The only vanishing mechanism}
% =====================================================================

In this stock there is one formula in the strip $0<\operatorname{Re}s<1$ that is allowed to be zero: the dual sum of the complementary pair.
The product cannot do it.
The circle cannot do it.
The dated table cannot do it.
The sum of the two ends is what remains.

\begin{theorem}[Only vanishing mechanism]
In the strip $0<\operatorname{Re}s<1$, the only formula in the stock that can vanish non-trivially is the dual sum of the pair.
\end{theorem}

A vanishing written as a pause-reading has one own-origin.
Write $\chi(s)$ for the ratio of local-circle samples at $1-s$ and at $s$.
The wrapping that makes the two sides of the pair the same sentence is the local circle together with the half-walk ratio $\pi$ and the factorial walk stretched off the integers.
Under that wrapping,
\[
\bigl|\chi(\sigma+it)\bigr|=1
\quad\text{if and only if}\quad
\sigma=\tfrac12.
\]
This does not wait on the next first-writing.
It is the section drawing, pulled onto the index.

\begin{figure}[ht]
\centering
\includegraphics[width=0.62\textwidth]{layer_e03_cut.png}
\caption{The cut is the even split of the unit.
HERE is the pause.
THERE is the elected lock.
Elevation looks at that hold from outside.}
\end{figure}

Off the lock, $|\chi|$ is the ratio of two radii.
Two own-origins.
A pause of this construction occupies one.
Therefore a pause-reading has $|\chi|=1$, hence $\sigma=\tfrac12$.

The off-lock dual sum, with $\chi$ compensating unequal ends so they can still add to zero, is not a pause of this construction.
It is two centres named at once.
Vacuum may name it.
Trace cannot walk it.
A vanishing is an arrival.
Privilege does not return by the back door.

\begin{theorem}[One centre]
A vanishing written as a pause-reading has $|\chi|=1$, hence $\sigma=\tfrac12$.
Combined with the weights lemma, the two ends meet only at half.
\end{theorem}

% =====================================================================
\section{Meetings sit on the cut}
% =====================================================================

\begin{lemma}[Meetings sit on the cut]
Every meeting of the derived unit that the construction can seat has real part $1/2$.
\end{lemma}

\begin{proof}
A meeting of the derived unit is a vanishing of a formula in the stock, from one HERE.
The only such formula in the strip is the dual sum of the pair.
A vanishing written as a pause-reading has one own-origin, hence $|\chi|=1$, hence $\sigma=\tfrac12$.

Suppose an off-lock pair $\{\rho,1-\rho\}$ with $\operatorname{Re}\rho\neq\tfrac12$.
The pairing $s\leftrightarrow 1-s$ names both ends.
Naming is vacuum.
A name is not a unit.
To seat the pair as a meeting the construction would have to occupy two own-origins at once.
The $n$-count does not.
Therefore the pair is not a seated meeting of this unit.
Therefore every meeting the derived unit can seat has real part $1/2$.
\end{proof}

\begin{theorem}[Midline]
Every meeting of the derived unit has real part $1/2$.
\end{theorem}

\begin{proof}
A meeting of the derived unit is a seated vanishing from one HERE.
Every meeting the construction can seat has real part $1/2$.
An unseated name is not a meeting of the unit.
\end{proof}

Dated crossings are unseated names.
They move when the station deepens.
They are not meetings.

% =====================================================================
\section{The reading that has a later name}
% =====================================================================

The three drawings, read together at finished stations, are one object.

Plan at finished stations is the unique-bag decomposition of the count.
Elevation at finished stations is the series~\eqref{eq:dirichlet}.
Where the series settles, the bags rewrite it as the product~\eqref{eq:euler}.

That object, on $\sigma>1$, is the series and the product already written.
The later filing of a name is not an ingredient of I, II, or III.
It is granted only after the drawings exist together.

Section supplies $\pi$ as half-walk over half-diameter, and the local circle.
The wrap $\chi$ is that section pulled onto the index.
Continuation through the half --- travel from a hard $s$ to $1-s$, where a writing arrives, and back --- is inversion, the first page of I, applied to the index.
The wrapped object satisfies
\[
X(s)=X(1-s).
\]
That is the lock, written on the two-part index.

Under that identification the dual sum in the strip is the writing the stock is allowed to set to zero.
A seated vanishing of that writing, from one HERE, with one centre, has real part $1/2$.

\begin{theorem}[The later name]
If the derived unit is identified with the classical limit obtained by reading plan, section, and elevation together through finished stations, then every non-trivial zero of that limit has real part $1/2$.
\end{theorem}

\begin{proof}
The identification names $Z(s)$ on $\sigma>1$ as the series already written, names the wrap $\chi$ as the factor that enforces $X(s)=X(1-s)$, and names a seated vanishing as a meeting of that continued object in the strip.
By the midline theorem every meeting of the derived unit has real part $1/2$.
The identification therefore carries the midline onto the later name.
\end{proof}

The identification is granted only after the drawings exist together.
It is not an ingredient of the construction.
It is a name for the reading.

A meeting off the cut would be two centres named as one arrival, or a second primary self-description of location.
The geometry does not admit either.

The first sentence is still the sentence.
Where is $1$?
On the cut.
Walk there and you are HERE.
The meetings that can be seated sit on that same line, because there is no other primary ratio for them to occupy.

% =====================================================================
\section{What a zero is, in this stock}
% =====================================================================

A zero of the later name is not a pebble in the disk.
It is not a last index.
It is not an extra room that appears when the count is declared finished.

It is a derived unit of measurement from one HERE: a pause, one own-origin, elevation looking at the closed unit from outside.
The formula that vanishes is the dual sum of the pair, wrapped so $|\chi|=1$.
The line it is allowed to occupy is the cut the unit already had.

That is all a seated vanishing is.
The word zeroing, used for hanging a live frame on the half-line, stays the hinge.
It is not this meeting.
The meeting is what the hinge permits after the three drawings have been read together.

% =====================================================================
\section{What stands}
% =====================================================================

Forced, and already forced in I--V, now with the elevation written:
look-back from a pause is $\log n$;
the weight of a tick is $n^{-s}$;
the series and the product are two drawings of one walk;
complementary interiors meet if and only if $\sigma=\tfrac12$;
a dated product never vanishes;
dated dual crossings are crop;
$|\chi|=1$ if and only if $\sigma=\tfrac12$;
a pause is one own-origin;
every seated meeting of the derived unit has real part $1/2$.

Named:
elevation, dated, dual sum, derived unit, meeting of the derived unit.

Granted as a reading, after the drawings exist together:
the identification of that derived unit with the classical limit of the same reading, and therefore the statement that every seated vanishing of that limit has real part $1/2$.

Not granted:
a second primary ratio;
two centres at one pause;
a name treated as a seat;
a dated crossing treated as a meeting;
a completed $n$-count as a first ingredient;
infinity as a last index;
the later name as an ingredient of I, II, or III.

The question is still the first sentence.
The answer is still the same unit.
Point to it.
You are pointing at the hold you already closed, at the cut that hold still admits, and at the elevation that looks at both from outside.
All three are $1$.
Arrive there and you are HERE.

\end{document}
